% Versuche V1 bis V11 nach Prüfebene und Umfang (Wiederholungen), nach Forschungsfragen gruppiert.
% Eigene Darstellung, keine Fremdquelle. Daten identisch mit tab:versuchsplan (Kapitel 5, Abschnitt sec:versuchsplan;
% Umfang: Vorschlag 16.11., plan-kap5.md Abschnitt 7.2). „offen“ = Umfang dort noch \pruefen (V3, V10).
% V4: je 10 Läufe zu 10 min; V6: 10 je Testfall; V9: 3 Durchläufe auf frischem System; V11: Einrichtung einmal,
% Ausführung Simulation 10, Roboter 3. Kann-Ebenen ohne Umfangsangabe in der Tabelle bleiben leer.
% Stil: ba-grafik.tex. Einbinden: % Versuche V1 bis V11 nach Prüfebene und Umfang (Wiederholungen), nach Forschungsfragen gruppiert.
% Eigene Darstellung, keine Fremdquelle. Daten identisch mit tab:versuchsplan (Kapitel 5, Abschnitt sec:versuchsplan;
% Umfang: Vorschlag 16.11., plan-kap5.md Abschnitt 7.2). „offen“ = Umfang dort noch \pruefen (V3, V10).
% V4: je 10 Läufe zu 10 min; V6: 10 je Testfall; V9: 3 Durchläufe auf frischem System; V11: Einrichtung einmal,
% Ausführung Simulation 10, Roboter 3. Kann-Ebenen ohne Umfangsangabe in der Tabelle bleiben leer.
% Stil: ba-grafik.tex. Einbinden: % Versuche V1 bis V11 nach Prüfebene und Umfang (Wiederholungen), nach Forschungsfragen gruppiert.
% Eigene Darstellung, keine Fremdquelle. Daten identisch mit tab:versuchsplan (Kapitel 5, Abschnitt sec:versuchsplan;
% Umfang: Vorschlag 16.11., plan-kap5.md Abschnitt 7.2). „offen“ = Umfang dort noch \pruefen (V3, V10).
% V4: je 10 Läufe zu 10 min; V6: 10 je Testfall; V9: 3 Durchläufe auf frischem System; V11: Einrichtung einmal,
% Ausführung Simulation 10, Roboter 3. Kann-Ebenen ohne Umfangsangabe in der Tabelle bleiben leer.
% Stil: ba-grafik.tex. Einbinden: % Versuche V1 bis V11 nach Prüfebene und Umfang (Wiederholungen), nach Forschungsfragen gruppiert.
% Eigene Darstellung, keine Fremdquelle. Daten identisch mit tab:versuchsplan (Kapitel 5, Abschnitt sec:versuchsplan;
% Umfang: Vorschlag 16.11., plan-kap5.md Abschnitt 7.2). „offen“ = Umfang dort noch \pruefen (V3, V10).
% V4: je 10 Läufe zu 10 min; V6: 10 je Testfall; V9: 3 Durchläufe auf frischem System; V11: Einrichtung einmal,
% Ausführung Simulation 10, Roboter 3. Kann-Ebenen ohne Umfangsangabe in der Tabelle bleiben leer.
% Stil: ba-grafik.tex. Einbinden: \input{Bilder/tikz/versuchsmatrix} in einer figure-Umgebung. Nicht skalieren.
\begin{tikzpicture}[ba]
  \tikzset{
    fest/.style={circle, draw=baGruen, fill=baFlaeche, line width=0.8pt, minimum size=6.4mm, inner sep=0pt, font=\footnotesize},
    kann/.style={circle, draw=baGrau, fill=white, line width=0.8pt, dash pattern=on 2pt off 1.5pt, minimum size=6.4mm, inner sep=0pt, font=\footnotesize},
    kopf/.style={font=\small, anchor=base},
    vlab/.style={font=\small, anchor=west},
    flab/.style={font=\small, anchor=west, text=baGrau},
    grenze/.style={draw=baHellgrau, ba normal},
  }
  % Spalten
  \def\xA{4.4}\def\xT{6.9}\def\xS{9.4}\def\xR{11.9}
  % Zeilen (y), Gruppen durch Abstand getrennt
  \def\ya{0}\def\yb{-0.75}\def\yc{-1.5}
  \def\yd{-2.55}\def\ye{-3.3}\def\yf{-4.05}\def\yg{-4.8}
  \def\yh{-5.85}\def\yi{-6.6}
  \def\yj{-7.65}\def\yk{-8.4}
  % ---------- Kopf ----------
  \node[kopf] at (\xA,0.55) {Analyse};
  \node[kopf] at (\xT,0.55) {Attrappe};
  \node[kopf] at (\xS,0.55) {Simulation};
  \node[kopf] at (\xR,0.55) {Roboter};
  \draw[ba pfeil, baGrau] ({\xT-0.9},1.45) -- ({\xR+0.9},1.45);
  \node[font=\footnotesize, text=baGrau, anchor=south] at (\xS,1.5) {Prüfebenen: wachsender Aufwand und wachsende Aussagekraft};
  % ---------- Gruppen F1 bis F4 ----------
  \foreach \y/\f/\s in {0/F1/sec:eval_f1, -2.55/F2/sec:eval_f2, -5.85/F3/sec:eval_f3, -7.65/F4/sec:eval_f4} {
    \node[flab] at (0,\y) {\f};
    \node[font=\footnotesize, text=baGrau, anchor=west] at (13.2,\y) {Abschnitt~\ref{\s}};
  }
  \foreach \y in {-2.0, -5.3, -7.1} \draw[grenze] (0,\y) -- (15.6,\y);
  \draw[grenze] (0,0.5) -- (15.6,0.5);
  \draw[grenze] (0,-8.85) -- (15.6,-8.85);
  % ---------- Versuchsnamen ----------
  \foreach \y/\v in {\ya/V1, \yb/V2, \yc/V3, \yd/V4, \ye/V5, \yf/V6, \yg/V7, \yh/V8, \yi/V9, \yj/V10, \yk/V11}
    \node[vlab] at (0.9,\y) {\v};
  % ---------- Markierungen ----------
  % F1
  \node[fest] at (\xA,\ya) {1};
  \node[fest] at (\xR,\yb) {3};
  \node[fest, font=\scriptsize] at (\xR,\yc) {offen};
  % F2
  \node[fest] at (\xT,\yd) {10};
  \node[fest] at (\xS,\yd) {10};
  \node[fest] at (\xS,\ye) {10};
  \node[kann] at (\xR,\ye) {3};
  \node[fest] at (\xS,\yf) {10};
  \node[fest] at (\xS,\yg) {10};
  \node[fest] at (\xR,\yg) {3};
  % F3
  \node[fest] at (\xT,\yh) {10};
  \node[kann] at (\xR,\yh) {3};
  \node[fest] at (\xS,\yi) {3};
  \node[kann] at (\xR,\yi) {};
  % F4
  \node[fest, font=\scriptsize] at (\xS,\yj) {offen};
  \node[kann] at (\xR,\yj) {};
  \node[fest] at (\xS,\yk) {10};
  \node[fest] at (\xR,\yk) {3};
  % ---------- Legende ----------
  \node[fest] (l1) at (0.35,-9.55) {};
  \node[font=\footnotesize, anchor=west] at ($(l1.east)+(0.1,0)$) {fest geplant};
  \node[kann] (l2) at (3.7,-9.55) {};
  \node[font=\footnotesize, anchor=west] at ($(l2.east)+(0.1,0)$) {Kann, abhängig von Zeit und Zugang zum Roboter};
  \node[font=\footnotesize, anchor=west, align=left] at (0.0,-10.45) {Zahl im Kreis: Wiederholungen bzw. Durchläufe; „offen“: Umfang noch festzulegen;\\leerer Kreis: Umfang nicht vorgegeben};
\end{tikzpicture}
 in einer figure-Umgebung. Nicht skalieren.
\begin{tikzpicture}[ba]
  \tikzset{
    fest/.style={circle, draw=baGruen, fill=baFlaeche, line width=0.8pt, minimum size=6.4mm, inner sep=0pt, font=\footnotesize},
    kann/.style={circle, draw=baGrau, fill=white, line width=0.8pt, dash pattern=on 2pt off 1.5pt, minimum size=6.4mm, inner sep=0pt, font=\footnotesize},
    kopf/.style={font=\small, anchor=base},
    vlab/.style={font=\small, anchor=west},
    flab/.style={font=\small, anchor=west, text=baGrau},
    grenze/.style={draw=baHellgrau, ba normal},
  }
  % Spalten
  \def\xA{4.4}\def\xT{6.9}\def\xS{9.4}\def\xR{11.9}
  % Zeilen (y), Gruppen durch Abstand getrennt
  \def\ya{0}\def\yb{-0.75}\def\yc{-1.5}
  \def\yd{-2.55}\def\ye{-3.3}\def\yf{-4.05}\def\yg{-4.8}
  \def\yh{-5.85}\def\yi{-6.6}
  \def\yj{-7.65}\def\yk{-8.4}
  % ---------- Kopf ----------
  \node[kopf] at (\xA,0.55) {Analyse};
  \node[kopf] at (\xT,0.55) {Attrappe};
  \node[kopf] at (\xS,0.55) {Simulation};
  \node[kopf] at (\xR,0.55) {Roboter};
  \draw[ba pfeil, baGrau] ({\xT-0.9},1.45) -- ({\xR+0.9},1.45);
  \node[font=\footnotesize, text=baGrau, anchor=south] at (\xS,1.5) {Prüfebenen: wachsender Aufwand und wachsende Aussagekraft};
  % ---------- Gruppen F1 bis F4 ----------
  \foreach \y/\f/\s in {0/F1/sec:eval_f1, -2.55/F2/sec:eval_f2, -5.85/F3/sec:eval_f3, -7.65/F4/sec:eval_f4} {
    \node[flab] at (0,\y) {\f};
    \node[font=\footnotesize, text=baGrau, anchor=west] at (13.2,\y) {Abschnitt~\ref{\s}};
  }
  \foreach \y in {-2.0, -5.3, -7.1} \draw[grenze] (0,\y) -- (15.6,\y);
  \draw[grenze] (0,0.5) -- (15.6,0.5);
  \draw[grenze] (0,-8.85) -- (15.6,-8.85);
  % ---------- Versuchsnamen ----------
  \foreach \y/\v in {\ya/V1, \yb/V2, \yc/V3, \yd/V4, \ye/V5, \yf/V6, \yg/V7, \yh/V8, \yi/V9, \yj/V10, \yk/V11}
    \node[vlab] at (0.9,\y) {\v};
  % ---------- Markierungen ----------
  % F1
  \node[fest] at (\xA,\ya) {1};
  \node[fest] at (\xR,\yb) {3};
  \node[fest, font=\scriptsize] at (\xR,\yc) {offen};
  % F2
  \node[fest] at (\xT,\yd) {10};
  \node[fest] at (\xS,\yd) {10};
  \node[fest] at (\xS,\ye) {10};
  \node[kann] at (\xR,\ye) {3};
  \node[fest] at (\xS,\yf) {10};
  \node[fest] at (\xS,\yg) {10};
  \node[fest] at (\xR,\yg) {3};
  % F3
  \node[fest] at (\xT,\yh) {10};
  \node[kann] at (\xR,\yh) {3};
  \node[fest] at (\xS,\yi) {3};
  \node[kann] at (\xR,\yi) {};
  % F4
  \node[fest, font=\scriptsize] at (\xS,\yj) {offen};
  \node[kann] at (\xR,\yj) {};
  \node[fest] at (\xS,\yk) {10};
  \node[fest] at (\xR,\yk) {3};
  % ---------- Legende ----------
  \node[fest] (l1) at (0.35,-9.55) {};
  \node[font=\footnotesize, anchor=west] at ($(l1.east)+(0.1,0)$) {fest geplant};
  \node[kann] (l2) at (3.7,-9.55) {};
  \node[font=\footnotesize, anchor=west] at ($(l2.east)+(0.1,0)$) {Kann, abhängig von Zeit und Zugang zum Roboter};
  \node[font=\footnotesize, anchor=west, align=left] at (0.0,-10.45) {Zahl im Kreis: Wiederholungen bzw. Durchläufe; „offen“: Umfang noch festzulegen;\\leerer Kreis: Umfang nicht vorgegeben};
\end{tikzpicture}
 in einer figure-Umgebung. Nicht skalieren.
\begin{tikzpicture}[ba]
  \tikzset{
    fest/.style={circle, draw=baGruen, fill=baFlaeche, line width=0.8pt, minimum size=6.4mm, inner sep=0pt, font=\footnotesize},
    kann/.style={circle, draw=baGrau, fill=white, line width=0.8pt, dash pattern=on 2pt off 1.5pt, minimum size=6.4mm, inner sep=0pt, font=\footnotesize},
    kopf/.style={font=\small, anchor=base},
    vlab/.style={font=\small, anchor=west},
    flab/.style={font=\small, anchor=west, text=baGrau},
    grenze/.style={draw=baHellgrau, ba normal},
  }
  % Spalten
  \def\xA{4.4}\def\xT{6.9}\def\xS{9.4}\def\xR{11.9}
  % Zeilen (y), Gruppen durch Abstand getrennt
  \def\ya{0}\def\yb{-0.75}\def\yc{-1.5}
  \def\yd{-2.55}\def\ye{-3.3}\def\yf{-4.05}\def\yg{-4.8}
  \def\yh{-5.85}\def\yi{-6.6}
  \def\yj{-7.65}\def\yk{-8.4}
  % ---------- Kopf ----------
  \node[kopf] at (\xA,0.55) {Analyse};
  \node[kopf] at (\xT,0.55) {Attrappe};
  \node[kopf] at (\xS,0.55) {Simulation};
  \node[kopf] at (\xR,0.55) {Roboter};
  \draw[ba pfeil, baGrau] ({\xT-0.9},1.45) -- ({\xR+0.9},1.45);
  \node[font=\footnotesize, text=baGrau, anchor=south] at (\xS,1.5) {Prüfebenen: wachsender Aufwand und wachsende Aussagekraft};
  % ---------- Gruppen F1 bis F4 ----------
  \foreach \y/\f/\s in {0/F1/sec:eval_f1, -2.55/F2/sec:eval_f2, -5.85/F3/sec:eval_f3, -7.65/F4/sec:eval_f4} {
    \node[flab] at (0,\y) {\f};
    \node[font=\footnotesize, text=baGrau, anchor=west] at (13.2,\y) {Abschnitt~\ref{\s}};
  }
  \foreach \y in {-2.0, -5.3, -7.1} \draw[grenze] (0,\y) -- (15.6,\y);
  \draw[grenze] (0,0.5) -- (15.6,0.5);
  \draw[grenze] (0,-8.85) -- (15.6,-8.85);
  % ---------- Versuchsnamen ----------
  \foreach \y/\v in {\ya/V1, \yb/V2, \yc/V3, \yd/V4, \ye/V5, \yf/V6, \yg/V7, \yh/V8, \yi/V9, \yj/V10, \yk/V11}
    \node[vlab] at (0.9,\y) {\v};
  % ---------- Markierungen ----------
  % F1
  \node[fest] at (\xA,\ya) {1};
  \node[fest] at (\xR,\yb) {3};
  \node[fest, font=\scriptsize] at (\xR,\yc) {offen};
  % F2
  \node[fest] at (\xT,\yd) {10};
  \node[fest] at (\xS,\yd) {10};
  \node[fest] at (\xS,\ye) {10};
  \node[kann] at (\xR,\ye) {3};
  \node[fest] at (\xS,\yf) {10};
  \node[fest] at (\xS,\yg) {10};
  \node[fest] at (\xR,\yg) {3};
  % F3
  \node[fest] at (\xT,\yh) {10};
  \node[kann] at (\xR,\yh) {3};
  \node[fest] at (\xS,\yi) {3};
  \node[kann] at (\xR,\yi) {};
  % F4
  \node[fest, font=\scriptsize] at (\xS,\yj) {offen};
  \node[kann] at (\xR,\yj) {};
  \node[fest] at (\xS,\yk) {10};
  \node[fest] at (\xR,\yk) {3};
  % ---------- Legende ----------
  \node[fest] (l1) at (0.35,-9.55) {};
  \node[font=\footnotesize, anchor=west] at ($(l1.east)+(0.1,0)$) {fest geplant};
  \node[kann] (l2) at (3.7,-9.55) {};
  \node[font=\footnotesize, anchor=west] at ($(l2.east)+(0.1,0)$) {Kann, abhängig von Zeit und Zugang zum Roboter};
  \node[font=\footnotesize, anchor=west, align=left] at (0.0,-10.45) {Zahl im Kreis: Wiederholungen bzw. Durchläufe; „offen“: Umfang noch festzulegen;\\leerer Kreis: Umfang nicht vorgegeben};
\end{tikzpicture}
 in einer figure-Umgebung. Nicht skalieren.
\begin{tikzpicture}[ba]
  \tikzset{
    fest/.style={circle, draw=baGruen, fill=baFlaeche, line width=0.8pt, minimum size=6.4mm, inner sep=0pt, font=\footnotesize},
    kann/.style={circle, draw=baGrau, fill=white, line width=0.8pt, dash pattern=on 2pt off 1.5pt, minimum size=6.4mm, inner sep=0pt, font=\footnotesize},
    kopf/.style={font=\small, anchor=base},
    vlab/.style={font=\small, anchor=west},
    flab/.style={font=\small, anchor=west, text=baGrau},
    grenze/.style={draw=baHellgrau, ba normal},
  }
  % Spalten
  \def\xA{4.4}\def\xT{6.9}\def\xS{9.4}\def\xR{11.9}
  % Zeilen (y), Gruppen durch Abstand getrennt
  \def\ya{0}\def\yb{-0.75}\def\yc{-1.5}
  \def\yd{-2.55}\def\ye{-3.3}\def\yf{-4.05}\def\yg{-4.8}
  \def\yh{-5.85}\def\yi{-6.6}
  \def\yj{-7.65}\def\yk{-8.4}
  % ---------- Kopf ----------
  \node[kopf] at (\xA,0.55) {Analyse};
  \node[kopf] at (\xT,0.55) {Attrappe};
  \node[kopf] at (\xS,0.55) {Simulation};
  \node[kopf] at (\xR,0.55) {Roboter};
  \draw[ba pfeil, baGrau] ({\xT-0.9},1.45) -- ({\xR+0.9},1.45);
  \node[font=\footnotesize, text=baGrau, anchor=south] at (\xS,1.5) {Prüfebenen: wachsender Aufwand und wachsende Aussagekraft};
  % ---------- Gruppen F1 bis F4 ----------
  \foreach \y/\f/\s in {0/F1/sec:eval_f1, -2.55/F2/sec:eval_f2, -5.85/F3/sec:eval_f3, -7.65/F4/sec:eval_f4} {
    \node[flab] at (0,\y) {\f};
    \node[font=\footnotesize, text=baGrau, anchor=west] at (13.2,\y) {Abschnitt~\ref{\s}};
  }
  \foreach \y in {-2.0, -5.3, -7.1} \draw[grenze] (0,\y) -- (15.6,\y);
  \draw[grenze] (0,0.5) -- (15.6,0.5);
  \draw[grenze] (0,-8.85) -- (15.6,-8.85);
  % ---------- Versuchsnamen ----------
  \foreach \y/\v in {\ya/V1, \yb/V2, \yc/V3, \yd/V4, \ye/V5, \yf/V6, \yg/V7, \yh/V8, \yi/V9, \yj/V10, \yk/V11}
    \node[vlab] at (0.9,\y) {\v};
  % ---------- Markierungen ----------
  % F1
  \node[fest] at (\xA,\ya) {1};
  \node[fest] at (\xR,\yb) {3};
  \node[fest, font=\scriptsize] at (\xR,\yc) {offen};
  % F2
  \node[fest] at (\xT,\yd) {10};
  \node[fest] at (\xS,\yd) {10};
  \node[fest] at (\xS,\ye) {10};
  \node[kann] at (\xR,\ye) {3};
  \node[fest] at (\xS,\yf) {10};
  \node[fest] at (\xS,\yg) {10};
  \node[fest] at (\xR,\yg) {3};
  % F3
  \node[fest] at (\xT,\yh) {10};
  \node[kann] at (\xR,\yh) {3};
  \node[fest] at (\xS,\yi) {3};
  \node[kann] at (\xR,\yi) {};
  % F4
  \node[fest, font=\scriptsize] at (\xS,\yj) {offen};
  \node[kann] at (\xR,\yj) {};
  \node[fest] at (\xS,\yk) {10};
  \node[fest] at (\xR,\yk) {3};
  % ---------- Legende ----------
  \node[fest] (l1) at (0.35,-9.55) {};
  \node[font=\footnotesize, anchor=west] at ($(l1.east)+(0.1,0)$) {fest geplant};
  \node[kann] (l2) at (3.7,-9.55) {};
  \node[font=\footnotesize, anchor=west] at ($(l2.east)+(0.1,0)$) {Kann, abhängig von Zeit und Zugang zum Roboter};
  \node[font=\footnotesize, anchor=west, align=left] at (0.0,-10.45) {Zahl im Kreis: Wiederholungen bzw. Durchläufe; „offen“: Umfang noch festzulegen;\\leerer Kreis: Umfang nicht vorgegeben};
\end{tikzpicture}
